% Passes 5816--5823: Fourier rank stratification of the q=5 matrix carrier.
\subsection{The $1+9+6$ line-carrier split is matrix-rank Fourier geometry}
\label{sec:pass5816-5823-matrix-fourier-rank}

The affine matrix model makes the full rational decomposition of the sixteen
Reye/Latin lines explicit.  The line set is the additive torsor
\[
 T=M_2(\mathbb F_2),
\]
and, using the Frobenius pairing
$\langle Y,M\rangle=\operatorname{tr}(Y^TM)$, define Walsh vectors
\[
 \boxed{
 v_Y=\sum_{M\in M_2(\mathbb F_2)}
 (-1)^{\langle Y,M\rangle}e_M.
 }
\]
They satisfy $\langle v_Y,v_Z\rangle=16\delta_{YZ}$.  The dual matrix labels
split by rank as
\[
 \boxed{16=1+9+6},
\]
namely one zero matrix, nine nonzero rank-one matrices, and six invertible
rank-two matrices.  This is exactly the carrier decomposition
\[
 \boxed{
 \mathbb Q^{16}_L
 =\mathbf1\oplus W_9^{\operatorname{rank}1}
 \oplus V_6^{\operatorname{rank}2}.
 }
\]

Indeed every nonzero rank-one label has the unique factorization
\[
 Y=\phi^Tw^T,
 \qquad0\ne\phi\in V^*,\quad0\ne w\in W,
\]
so the nine labels form a literal $3\times3$ grid.  Define the point- and
heavy-fibre Walsh vectors
\[
 u_{w,\phi}=\sum_{x\in V}(-1)^{\phi(x)}e_{(w,x)},
 \qquad
 h_{w,\phi}=\sum_{\psi\in W^*}(-1)^{\psi(w)}e_{(\phi,\psi)}.
\]
For point--line Reye incidence $R$, point--heavy incidence $H$, and
line--heavy disjointness incidence $D$, exact summation gives
\[
 \boxed{R^Tu_{w,\phi}=v_{\phi^Tw^T}},
 \qquad
 \boxed{H^Tu_{w,\phi}=-2h_{w,\phi}},
 \qquad
 \boxed{Dh_{w,\phi}=v_{\phi^Tw^T}}.
\]
Thus the point, heavy and line realizations of the common $W_9$ are exactly the
same nine rank-one additive characters, transported by finite Radon transforms.

The six invertible labels give the complementary line-only sector.  For every
$Y\in GL_2(2)$,
\[
 \boxed{Rv_Y=0},
 \qquad
 \boxed{D^Tv_Y=0},
\]
while no nonzero rank-one Walsh vector is killed by either transform.  Hence
\[
 \boxed{
 V_6=\operatorname{span}\{v_Y:\det Y=1\}
 =\ker R=\ker D^T.
 }
\]

The full affine carrier also becomes monomial in this basis.  For
$g=(X,A,B)$,
\[
 \boxed{
 gv_Y=(-1)^{\langle Y',X\rangle}v_{Y'},
 \qquad Y'=A^{-T}YB^T.
 }
\]
This identity was checked on all $576\cdot16=9216$ element/label pairs.  The
rank-one and rank-two sector characters satisfy
\[
 \boxed{\langle\chi_9,\chi_9\rangle=1},
 \qquad
 \boxed{\langle\chi_6,\chi_6\rangle=1},
 \qquad
 \boxed{\langle\chi_9,\chi_6\rangle=0}.
\]
Because both sectors are explicitly realized over $\mathbb Q$ by integer
Walsh vectors, each is absolutely irreducible.

Restriction to the normal translation group $T=M_2(\mathbb F_2)_+$ is
especially transparent:
\[
 \boxed{
 W_9\downarrow_T
 =\bigoplus_{\operatorname{rank}Y=1}\chi_Y,
 }
 \qquad
 \boxed{
 V_6\downarrow_T
 =\bigoplus_{\det Y=1}\chi_Y.
 }
\]
Thus the normal sixteen-group itself spectrally distinguishes the common-nine
and line-only-six sectors.

Transpose acts on the Fourier labels by
\[
 \boxed{\Theta v_Y=v_{Y^T}}.
\]
For $Y=\phi^Tw^T$ this sends $Y$ to $w^T\phi^T$, so the same outer involution
that exchanges point and heavy carriers swaps the two factors of the
$3\times3$ rank-one grid and preserves the six invertible labels as a set.

Finally, over $\mathbb F_2$ each nonzero matrix is already a unique projective
point of $\mathrm{PG}(3,2)$.  The determinant-zero locus consists of precisely
the nine rank-one points and its complement of the six invertible points:
\[
 \boxed{15=9+6}.
\]
The nine-point locus has the explicit product parametrization
$(\phi,w)\in3\times3$, the standard hyperbolic-quadric/grid rank geometry of
$2\times2$ matrices.

\paragraph{Prior art and evidence boundary.}
Finite translation/Fourier rank schemes are classical; an explicit modern
reference for finite bilinear and quadratic-form translation schemes is
K.-U.~Schmidt, arXiv:1803.04274.  The $M_2(\mathbb F_2)$ two-qubit
projective-line use is also prior art in Saniga--Planat--Pracna,
arXiv:quant-ph/0611063.  The new repository statement is the exact
identification of the already-certified q=5 Reye/heavy carrier modules and
incidence maps with these rank-stratified Walsh sectors.  The dimensions
$1+9+6$ are not particle multiplets, qubit states or physical sectors absent a
separately verified physical map.
